\chapter{Спектральный нейронный оператор Фанаскова и Оселедца}
\label{sno:chapter}

Пусть оператор $\mathcal G$ каждой входной функции $f$ сопоставляет
выходную функцию $g$. В спектральном нейронном операторе
(\emph{Spectral Neural Operator}, SNO) Фанаскова и Оселедца обе функции
представлены конечными рядами. Чтобы приближать $\mathcal G$,
достаточно научиться по коэффициентам первого ряда определять
коэффициенты второго. Авторы рассматривают полиномы Чебышёва и ряды
Фурье. Разберём сначала полиномиальный случай.
Общее определение дано в
\href{https://arxiv.org/html/2205.10573v2\#S4}{разделе~4 статьи, формула (3)},
а устройство сети --- в приложении~C.1.3, формулы (19)--(22).

\section{Коэффициенты входной и выходной функций}
\label{sno:bases}

На отрезке $[-1,1]$ положим $T_j(\cos t)=\cos(jt)$.
Первые три полинома Чебышёва равны $T_0=1$, $T_1=x$ и $T_2=2x^2-1$.
Условимся записывать ряд в виде $f_p=\sum_{j=0}^p d_j T_j$.
Тогда, например,
\[
 f(x)=1+x+x^2=\tfrac32T_0(x)+T_1(x)+\tfrac12T_2(x),
 \qquad d=(3/2,1,1/2)^\top .
\]
При $x=0$ имеем $T_2(0)=-1$, и ряд даёт исходное значение $f(0)=1$.

Ортогональность полиномов с весом $\rho(x)=(1-x^2)^{-1/2}$
определяет способ вычисления $d$. Если
$h_j=\int_{-1}^1 T_j(x)^2\rho(x)\,dx$, то
\begin{equation}
 d_j=\frac{1}{h_j}\int_{-1}^1 f(x)T_j(x)\rho(x)\,dx,
 \qquad h_0=\pi,\quad h_j=\pi/2\ (j\ge1).
 \label{sno:cheb-coefficients}
\end{equation}
Для полинома степени не выше $p\ge1$ те же числа восстанавливаются
по значениям в $p+1$ узлах $x_m=\cos(\pi m/p)$.
Например, при $p=2$ значения $f(1)=3$, $f(0)=1$, $f(-1)=1$ дают
$d_0+d_1+d_2=3$, $d_0-d_2=1$ и $d_0-d_1+d_2=1$.
Из первой и третьей строк получаем $d_1=1$, а из двух оставшихся
условий $d_0+d_2=2$ и $d_0-d_2=1$ находим $d_0=3/2$ и $d_2=1/2$.
Для произвольной непрерывной функции значения в этих узлах определяют
интерполирующий полином. Его коэффициенты могут отличаться от интегралов
\eqref{sno:cheb-coefficients}.

Обозначим извлечение $p+1$ входных коэффициентов через $E_p$, а
вычисление выходного ряда через $D_r$. Сеть $N_\vartheta$
определяет приближённый оператор
\begin{equation}
 f\ \xrightarrow{E_p}\ d\in\R^{p+1} \xrightarrow{N_\vartheta}\ \widehat b\in\R^{r+1} \xrightarrow{D_r}\ \widehat g(x)=\sum_{i=0}^{r}\widehat b_i T_i(x).
 \label{sno:operator-factorization}
\end{equation}
Степени $p$ и $r$ выбираются отдельно. При обучении известны пары
$(f,g)$, следовательно, доступны входные и эталонные выходные
коэффициенты. По этим парам подбирают общие веса $\vartheta$.
Для новой функции $f$ веса уже зафиксированы, а $\widehat b$
вычисляется заново. После этого $\widehat g$ можно вычислить в любой
точке отрезка. Число степеней свободы ряда при сгущении сетки остаётся
равным $r+1$.

Пусть, например, требуется учиться дифференцированию полиномов степени
не выше двух. Производные первых трёх мод равны $0$, $T_0$ и $4T_1$,
поэтому эталонное отображение коэффициентов в этой задаче известно:
\[
 b=\begin{pmatrix}0&1&0\\0&0&4\end{pmatrix}d .
\]
Входу $(3/2,1,1/2)^\top$ соответствует метка $(1,2)^\top$.
Её декодирование даёт $g(x)=1+2x$, то есть производную исходного
полинома. На такой паре можно проверить сразу оба преобразования
и смысл регрессии: для этого входа сеть должна приблизить числа 1 и 2.

Для периодической функции периода~2 ту же схему задают
$g_\nu(x)=e^{\ii\pi\nu x}$, $|\nu|\le K$. В этой нормировке
$d_\nu=\tfrac12\int_{-1}^1f(x)e^{-\ii\pi\nu x}\,dx$.
Для тригонометрического полинома с указанной полосой интеграл можно
заменить дискретным средним. Если $J=2K+1$ и $x_m=-1+2m/J$, то
\[
 d_\nu=\frac1J\sum_{m=0}^{J-1}f(x_m)e^{-\ii\pi\nu x_m}.
\]
Например, у $f(x)=\cos(\pi x)$ коэффициенты $d_{-1}=d_1=1/2$,
а остальные равны нулю. Это частный случай симметрии
$d_{-\nu}=\overline{d_\nu}$ для вещественных функций. Авторский код
пользуется этой симметрией и хранит неотрицательные частоты
\href{https://github.com/VLSF/SNO/blob/4362226284f22f6f456d429ec4ddec5b8596e894/functions/Fourier.py\#L43-L56}{(прямое и обратное преобразования)}.

\section{Как один слой действует на функции}
\label{sno:mixed-layer}

Пусть на входе слоя $\ell$ функций, каждая представлена $k$ модами.
Запишем их коэффициенты в матрицу $U\in\R^{k\times\ell}$.
Столбец задаёт одну функцию, а строка собирает вклад одной моды во
всех функциях. Внутри сети такие функции называют скрытыми каналами.
Разные примеры обучающей выборки имеют отдельные матрицы $U$.

Исходная сеть состоит из трёх подсетей $N_3\circ N_2\circ N_1$.
\[
 \underbrace{U}_{\text{входные каналы}}
 \xrightarrow{N_1}
 \underbrace{V}_{\text{скрытые каналы}}
 \xrightarrow{N_2}
 \underbrace{Y}_{\text{новые моды и каналы}}
 \xrightarrow{N_3}
 \underbrace{\widehat U}_{\text{выходные каналы}} .
\]
В $N_1$ и $N_3$ каналы смешиваются одинаково для каждой строки:
$U^+=\sigma(UC+\bm 1 a^\top)$, где
$C\in\R^{\ell\times m}$, $a\in\R^m$, $\bm 1\in\R^k$.
Такое смещение сдвигает каждый коэффициент данного канала на одно и то же число.
Например, прибавление $a$ к обеим модам $T_0,T_1$ добавляет к
восстановленной функции $a+ax$. В $N_2$ добавляется матрица $B$,
связывающая разные моды. Типовой
внутренний слой статьи имеет вид
\begin{equation}
 U^+=\sigma(BUC+\beta),\qquad
 B\in\R^{s\times k},\quad
 C\in\R^{\ell\times m},\quad
 \beta\in\R^{s\times m}.
 \label{sno:paper22}
\end{equation}
Здесь $C$ обозначает матрицу $A$ из формулы (22) статьи.
Умножение справа меняет число каналов с $\ell$ на $m$,
умножение слева --- число мод с $k$ на $s$.
Каждое выходное число до активации равно
\[
 z_{i\gamma}=\sum_{j=1}^k\sum_{\alpha=1}^{\ell}
         B_{ij}U_{j\alpha}C_{\alpha\gamma}+\beta_{i\gamma}.
\]
Тем самым один слой задаёт веса вида $B_{ij}C_{\alpha\gamma}$.
При записи $U$ столбец за столбцом его линейная часть имеет матрицу
$C^\top\otimes B$, составленную из блоков $C_{\alpha\gamma}B$.
Слой содержит $sk+\ell m+sm$ обучаемых вещественных чисел, включая
смещение. Аффинное отображение с отдельным весом для каждой пары
входной и выходной координат потребовало бы $(k\ell)(sm)+sm$ чисел.

\clearpage
\textbf{Пример.}
Пусть базис состоит из $T_0,T_1$, а два входных канала представляют
функции $1$ и $2x$. Выберем для учебного расчёта
\[
 U=\begin{pmatrix}1&0\\0&2\end{pmatrix},\quad
 C=\begin{pmatrix}1&-1\\1&1\end{pmatrix},\quad
 B=\begin{pmatrix}1&1\\1&-1\end{pmatrix},\quad
 \beta=\begin{pmatrix}0&0\\1&0\end{pmatrix}.
\]
После умножения справа и смешивания мод получаем
\[
 UC=\begin{pmatrix}1&-1\\2&2\end{pmatrix},
 \qquad BUC=\begin{pmatrix}3&1\\-1&-3\end{pmatrix},
 \qquad Z=BUC+\beta=\begin{pmatrix}3&1\\0&-3\end{pmatrix}.
\]
Применим поэлементно $s(t)=\log(1+e^t)$ (softplus). Тогда
\[
 U^+=s(Z)\approx\begin{pmatrix}3{,}04859&1{,}31326\\0{,}69315&0{,}04859\end{pmatrix}.
\]
Первый выходной канал есть $3{,}04859+0{,}69315x$. В точке $x=1$
он принимает значение $3{,}74174$. Число $0{,}69315$ в нижней строке
определяет наклон этой функции. Так по одной записи $U^+$
можно отдельно проверить вычисление слоя и восстановление функции.

В авторской программе последний слой \emph{каждой} подсети возвращает
аффинный результат $Z$ без активации. Для комплексных каналов внутри подсети
используется $s(\operatorname{Re}z)+\ii s(\operatorname{Im}z)$.
При обучении усредняют нормы Фробениуса ошибок таблиц по примерам.
Вычисление нормы состоит из суммирования квадратов модулей ошибок
коэффициентов и извлечения квадратного корня. Для отчётных метрик
статьи функции восстанавливают на общей равномерной сетке
(приложение~C.3).
\href{https://github.com/VLSF/SNO/blob/4362226284f22f6f456d429ec4ddec5b8596e894/architectures/SNO_1D.py}{Вещественный}
и
\href{https://github.com/VLSF/SNO/blob/4362226284f22f6f456d429ec4ddec5b8596e894/architectures/fSNO_1D.py}{комплексный код для случая 1D}.

\section{Связь матрицы мод с интегральным оператором}

Почему обычная матрица $B$ описывает оператор, действующий на функции?
Рассмотрим один канал в базисе Чебышёва, с нумерацией мод от нуля.
Линейное преобразование $d\mapsto Bd$ определяет функцию
$(\mathcal B f)(x)=\sum_{i,j} T_i(x)B_{ij}d_j$.
Подставив интеграл \eqref{sno:cheb-coefficients} и вынеся конечную
сумму под знак интеграла, получим
\begin{equation}
 (\mathcal B f)(x)=\int_{-1}^1 K_B(x,y)f(y)\rho(y)\,dy,
 \qquad
 K_B(x,y)=\sum_{i,j} T_i(x)B_{ij}\frac{T_j(y)}{h_j}.
 \label{sno:kernel}
\end{equation}
Таким образом, элементы матрицы задают коэффициенты ядра в двух
пространственных переменных. Образ $\mathcal B$ лежит в линейной
оболочке выходных полиномов. Ограничение числа мод одновременно
ограничивает ранг оператора.

Положение нелинейности тоже имеет значение. Возведение функции
$f(x)=\cos(\pi Kx)$ в квадрат порождает гармонику $2K$, поскольку
$f^2=(1+\cos(2\pi Kx))/2$. Возьмём $J=2K+1$ и узлы $x_m=2m/J$,
$m=0,\ldots,J-1$, на одном периоде. В этих узлах
$e^{\ii\pi(2K)x_m}=e^{-\ii\pi x_m}$, поскольку множитель
$e^{\ii \pi J x_m}$ равен 1. Поэтому высокую гармонику невозможно
отличить по отсчётам от гармоники с номером $-1$. Это наложение
частот (\emph{aliasing}) при дискретизации. В слое~\eqref{sno:paper22}
нелинейность применяется непосредственно к числам таблицы. После
неё ряд имеет ровно столько позиций, сколько строк в $U^+$.
Если истинный ответ содержит отброшенные моды, ошибка усечения
сохраняется даже при точном прогнозе всех оставленных коэффициентов.

В SVD-SNO v2 моды строятся по обучающим полям методом POD.
В блоке появляются локальное преобразование каналов и остаточная
связь (глава~\ref{code:implementation}).
